% from main.tex
%%%%%%%%%%  scan idx 136  |  folio 123  %%%%%%%%%%
\textbf{Proposition 4.10} \emph{La catégorie \(\Parf(\mathcal C,\mathcal C_0)\) est une sous-\(S\)-catégorie triangulée de \(\Dpx(\mathcal C)\) (1.1.4). Le foncteur canonique}
\[
  \Kpx^b(\mathcal C_0)\longrightarrow\Parf(\mathcal C,\mathcal C_0)
\]
(resp. \(\Tr(\Kpx^b(\mathcal C_0))\to\Tr(\Parf(\mathcal C,\mathcal C_0))\) (1.1.3)) \emph{induit une \(S\)-équivalence sur les champs associés ([G] II 2.2.4), ce qui signifie} :
\begin{enumerate}[label=(\roman*)]
\item \emph{pour tout \(X\in\Ob S\) et tous} \(E,F\in\Ob\Kpx^b(\mathcal C_{0X})\) (resp. \(\Tr(\Kpx^b(\mathcal C_{0X}))\)), \emph{la flèche canonique}
\[
  \underline{\Hom}_{\Kpx^b(\mathcal C_0)}(E,F)
  \longrightarrow
  \underline{\Hom}_{\Dpx(\mathcal C)}(E,F)
\]
(\emph{resp.} \(\underline{\Hom}_{\Tr(\Kpx^b(\mathcal C_0))}(E,F)\to\underline{\Hom}_{\Tr(\Dpx(\mathcal C))}(E,F)\)) \emph{est un isomorphisme, et}
\item \emph{tout objet de \(\Parf(\mathcal C,\mathcal C_0)\) (resp. \(\Tr(\Parf(\mathcal C,\mathcal C_0))\)) est localement isomorphe dans \(\Dpx(\mathcal C)\) (resp.} \(\Tr(\Dpx(\mathcal C))\) \emph{à un objet de \(\Kpx^b(\mathcal C_0)\) (resp.} \(\Tr(\Kpx^b(\mathcal C_0))\).
\end{enumerate}

\emph{Preuve.} L'assertion (i) non respée est un cas particulier de (4.6) et l'assertion (ii) non respée est vraie par définition ((4.7) et (4.9)). Il est clair d'autre part que \(\Parf(\mathcal C,\mathcal C_0)\) est stable par translation des degrés. Pour prouver que \(\Parf(\mathcal C,\mathcal C_0)\) est une sous-\(S\)-catégorie triangulée de \(\Dpx(\mathcal C)\), il suffit donc de montrer que si
\[
\begin{tikzcd}
& G \arrow[dl,"+1"'] & \\
E \arrow[rr,"u"] && F \arrow[ul]
\end{tikzcd}
\tag{*}
\]
est un triangle distingué de \(\Dpx(\mathcal C)\) tel que \(E\) et \(F\) soient parfaits, alors \(G\) est
